\scanpage{271}

équivalence \inserted{\unsure{de} 2-cat,} i.e. qu'il est \unsure{3} fidèle.

Bien entendu on est ramené quand \unsure{\(\overline{\mathrm m}\)} au cas de groupoïdes \inserted{\unsure{« connexes »}} \(\Pi_\Gamma\) définis par un groupe \inserted{(\unsure{fini} s'il le faut)}, et on est aussi ramené par des arguments \unsure{cat.} triviaux à regarder le cas des diagrammes \(B_{\unsure{D^*}}\to B_U\) avec \(B_U\) connexe.

Du côté géométrique, la situation \unsure{serait} donnée par un \(U=X\setminus S\) de type \((g,\nu)\) (\((g,\nu)\ne(0,0)\)), et une opération de \(\Gamma\) dessus --- OPS \(U=U_{g,\nu}\), donc on donne \(\Gamma\to A_{g,\nu}\). Pour la question des \unsure{\(i\)-fidélité} avec \(i\leq 2\), OPS qu'il s'agit du \unsure{\(\overline{\mathrm m}\)} groupe \(\Gamma\) qui opère\unsure{;} \deleted{\illegible{4}}

\unsure{1)} \uline{0-fidélité}. \deleted{\illegible{1}} Soient \(U,U'\) avec op. de \(\Gamma\) dessus, \(f,g:U\stackrel{\sim}{\longrightarrow}U'\) commutant à \(\Gamma\), et \(\alpha,\beta\) deux homotopies \inserted{\unsure{de \(f\) à \(g\)}} i.e. deux chemins de \(f\) à \(g\) dans \(\uline{\mathrm{Isom}}_\Gamma(U,U')\), on suppose que dans les descriptions \unsure{topossiques}
\diagram{\begin{tikzcd}[ampersand replacement=\&,column sep=5.5em,row sep=3.5em]
B_{D^*,\Gamma}
 \arrow[r]
 \arrow[d,bend right=16,"f_{D^*,\Gamma}"']
 \arrow[d,bend left=16,"g_{D^*,\Gamma}"]
\& B_{U,\Gamma}
 \arrow[r,no head]
 \arrow[d,bend right=16,"f_{U,\Gamma}"']
 \arrow[d,bend left=16,"g_{U,\Gamma}"]
\& B_\Gamma \arrow[d,equal]
\\
\unsure{B_{D'^*,\Gamma}} \arrow[r,no head]
\& B_{U',\Gamma} \arrow[r,no head]
\& B_\Gamma
\end{tikzcd}}
on a \(\alpha_*^{D^*}=\beta_*^{D^*}:f_{D^*}\to g_{D^*}\) et \(\alpha_*^U=\beta_*^U:f_U\to g_U\). A prouver que \(\alpha\) et \(\beta\) sont homotopes.

\margin{270}

\scanpage{272}
\mspage{137}

\margin{i.e. \(f\in A_{g,\nu}\) commutant à l'image de \(\Gamma\) par \(\Gamma\to A_{g,\nu}\).}
OPS \(g=\mathrm{id}\), \(U=U_{g,\nu}\), donc \(f\in A_{g,\nu}^{\Gamma}\), et on a deux chemins \(\alpha,\beta\) de \(1\) à \deleted{\illegible{1}} \(f\) dans l'espace \(A_{g,\nu}^{\Gamma}\). On suppose que les deux isomorphismes correspondants entre identité \inserted{de \unsure{\(\Pi_{D^*}\)}} et \unsure{\(f_{D^*,\Gamma}\)} : \(\Pi_{\deleted{\unsure{B}}\unsure{D^*_{S_\nu},\Gamma}}\to\Pi_{\unsure{D^*_{S_\nu},\Gamma}}\) d'une part, entre identité de \(\Pi_{U,\Gamma}\) et \(f_{U,\Gamma}\) d'autre part, sont les \(\overline{\mathrm m}\), on veut prouver que \(\alpha\) et \(\beta\) sont homotopes. Bien sûr, l'hyp. \deleted{\unsure{sur \(\alpha,\beta\)}} relative à l'action de \(\Gamma\) est vérifiée a fortiori en se restreignant à un groupe plus petit, p.ex. \deleted{\unsure{\(\Gamma\)}} \(\Gamma'=1\), et en fait on voit qu'elle est équivalente pour \(\Gamma\), et pour \unsure{son sous-groupe \(1\)}. \deleted{\illegible{1}} le résultat déjà connu (\unsure{la} \illegible{1}) pour \(\Gamma=1\), montre que ceci signifie que si deux chemins \inserted{dans \unsure{\(A_{g,\nu}^{\Gamma}\)}} de \(1\) à \(f\) sont homotopes dans \(A\), ils le sont dans \(A^{\Gamma}\) ou encore que
\margin{\unsure{posons} \(A=A_{g,\nu}\)}
\[
\fbox{\(\pi_1(A^{\Gamma})\to\pi_1(A)\ \mbox{est injectif.}\)}
\]
Dans le cas \(\pi_1(A)=0\) (cas \((g,\nu)\) \unsure{anabélien}) cela revient donc à prouver que \(\pi_1(A^{\Gamma})=0\).

\margin{271}

\scanpage{273}

2) \uline{1-fidélité}. Cela signifie \inserted{(en plus de la 0-fid)} que \unsure{\(\forall\)} isomorphisme entre \deleted{\unsure{\(f\)}} \(B_f\) et \(B_g\) provient d'un \deleted{\unsure{chem}} chemin de \(f\) à \(g\). Avec les réductions précédentes, OPS \(g=\mathrm{id}\), donc on a \(f\in A^{\Gamma}\), \deleted{\unsure{et l'on}} \unsure{\(f^*:(B_{D^*}\to B_U)\to(B_{D^*}\to B_U)\)} (\unsure{commutant} \deleted{\unsure{à}} \(\Gamma\) \inserted{\unsure{respectant \(\Gamma\)}}), et on a un iso. avec l'identité --- sauf erreur cette condition

\margin{272}

\scanpage{274}
\mspage{138}

Soit \(A\) un groupe topologique, d'où deux invariants
\[
\pi_0=\pi_0(A),\qquad \pi_1=\pi_1(A)
\]
le premier un groupe \inserted{pas néc. comm.}, \deleted{\illegible{1}} le deuxième un groupe commutatif sur lequel \(\pi_0\) opère. On définit une Gr-catégorie \inserted{\(\uline{A}\)} d'invariants \(\pi_0,\pi_1\) et correspondant à cette opération de \(\pi_0\) sur \(\pi_1\), en prenant comme catégorie sous-jacente le \deleted{\illegible{3}} groupoïde \deleted{\unsure{Hom}} fondamental (naïf) de \(A\), et comme foncteur de composition celui déduit de \(A\times A\to A\). (l'associativité dans \(\uline{A}=\Pi_1 A\) est stricte \ldots)

Ceci posé, tout torseur sur \(A\) définit un 1-torseur sur la gr-catégorie \(\uline{A}\). Plus généralement, pour \(\forall\) espace topologique \(S\) (ou \(\forall\) topos qui est loc. un espace topologique) tout torseur sur \(S\) de groupe \(A_S\) définit un \inserted{\(\uline{A}_S\)-}champ sur \(S\) qui est un \deleted{\unsure{\(\uline{A}_S\)}} champ en \(A_S\)-torseurs. Si \(\pi_i(A)=0\) pour \(i\geq2\), alors on trouve

\margin{273}

\scanpage{275}

ainsi \inserted{pour \unsure{\(\forall n\in N,\ n\geq2\)}} une \deleted{\illegible{1}} \uline{équivalence} entre la \unsure{\(n\)}-catégorie des torseurs sur \(S\) de groupe \(A_S\), et la \unsure{\(n\)}-catégorie déduite de la 2-catégorie des \inserted{\unsure{gerbes}} \(\uline{A}_S\)-torseurs en \unsure{la} prolongeant de façon discrète \ldots

Mais considérons \deleted{\illegible{1}} un groupe \(\Gamma\), et considérons la classification des \(A_{B_\Gamma}\)-torseurs sur le topos classifiant --- i.e. celle des \(A\) torseurs, avec une opération de \(\Gamma\) dessus (commutant à l'action de \(A\)). Ces objets forment une 2-catégorie, dont les composantes connexes correspondent aux classes de conjugaison d'homomorphismes de \(\Gamma\) dans \(A\). Si on a \deleted{deux} \inserted{un} hom \(\Gamma\to A\), il définit \deleted{une} une action sur le torseur trivial. \(1_A\) Si \(u,v\) sont deux hom, les isom de \((1_A,u)\) sur \((1_A,v)\) correspondent aux éléments de l'ens. \unsure{«}
\[
\unsure{\mathrm{Transp}_A}(u,v)=\{g\in A\mid v=\mathrm{int}\,g.u\}
\]
Mais on \unsure{fait|fera} une catégorie 0-isotopique,

\margin{274}

\scanpage{276}
\mspage{139}

en remplaçant \(\mathrm{Transp}(u,v)\) par son \(\pi_0\), et une catégorie 1-isotopique en \inserted{le} remplaçant par son groupoïde fondamental \(\Pi_1\mathrm{Transp}(u,v)\). Donc un isom de \(u\) à \(v\) est encore un pt de \(\mathrm{Transp}(u,v)\), mais si on a deux tels isomorphismes \inserted{\(f,g\)}, les isom \(f\simeq g\) sont les classes de \unsure{chemins} dans \(\mathrm{Transp}(u,v)\) de \(f\) à \(g\). Ainsi, la catégorie \(\uline{\mathrm{Aut}}(u)\) est équivalente à \(\Pi_1\underbrace{\unsure{\mathrm{Trans}}(u,u)}_{A^u}\) \deleted{est}
\[
\uline{\mathrm{Aut}}(u)\simeq\Pi_1\unsure{\mathrm{Trans}}(u)
\]
et la \deleted{\illegible{1}} catégorie \(\mathrm{Isom}(u,v)\) est soit vide (si \(u,v\) ne sont conjugués), soit un torseur sous \(\uline{\mathrm{Aut}}(u)\).

Mais pour tout torseur \inserted{\(P\)} sur \(A\) sur lequel \(\Gamma\) opère, \unsure{le} \deleted{\unsure{associé un}} 1-torseur \deleted{sous} \(\Pi_1P\) sur \(\Pi_1A\) \inserted{\(=\uline{A}\)} est muni d'opérations de \(\Gamma\), d'où un

\margin{275}

\scanpage{277}

\unsure{\((\Gamma,-\uline{A})\)} torseur. On trouve ainsi un foncteur de 2-catégories
\diagram{\begin{tikzcd}[ampersand replacement=\&,column sep=large]
{\begin{array}{l}
\mbox{2 catégorie des }\deleted{\mbox{torseurs}}\\
\Gamma-A\mbox{-torseurs}\\
\mbox{[ou encore, des hom }\Gamma\to A\mbox{]}
\end{array}}
\arrow[r] \&
{\begin{array}{l}
\mbox{2 catégorie des 1-torseurs sous }\unsure{\uline{A}},\\
\mbox{avec opérations de }\Gamma\mbox{ dessus.}
\end{array}}
\end{tikzcd}}

En somme, on vient de répéter sur le topos classifiant \(B_\Gamma\) la construction faite \inserted{plus haut} pour un espace topologique \(S\). On aimerait avoir \inserted{\unsure{exprimées}} des conditions pour que cet \deleted{hom} \inserted{2-foncteur} soit une 2-équivalence. Pour ceci, il conviendrait d'abord d'avoir une compréhension de la classification des classes \unsure{d'équivalence|d'équivalences} d'objets de la deuxième catégorie --- qui sont, eux, de nature purement algébrique, en termes de la Gr-catégorie \(\uline{A}\) et de \(\Gamma\). \unsure{Quitte à} se borner à des

\margin{276}

\scanpage{278}
\mspage{140}

torseurs triviaux sur \(\uline{A}\) (ce qui est \unsure{licite}) il faudrait expliciter ce que signifie que \(\Gamma\) opère sur le torseur trivial. On constate que cela signifie qu'on a un hom \inserted{de Gr-cat.} de \inserted{la Gr-cat. discrète définie par} \(\Gamma\) dans \(\uline{A}\). Donc la 2-catégorie envisagée est celle dont les objets sont les hom de Gr-catégories
\diagram{\begin{tikzcd}[ampersand replacement=\&]
\Gamma \arrow[r,"u","{\deleted{\unsure{\underset{v}{\longrightarrow}}}}"'] \& \uline{A}
\end{tikzcd}}
et pour deux tels morphismes \(u,v\), il faut définir (non seulement un ensemble \(\mathrm{Hom}(u,v)\) mais) une catégorie \(\uline{\mathrm{Hom}}(u,v)\), comme la sous-catégorie \(\uline{\mathrm{Transp}}(u,v)\) de \(\uline{A}\), à définir ad-hoc.

Considérons pour simplifier le cas où \(\pi_1=0\), alors \(\uline{A}\) se réduit à \(\pi_0\), et la cat. des hom \(\Gamma\to\uline{A}\) à celle des hom \(\Gamma\to\pi_0\). Le

\margin{277}

\scanpage{279}

fait que le foncteur de 2-catégories plus haut \deleted{\unsure{qu'il}} soit une équivalence de \unsure{2|1}-catégories se traduit alors peu à peu ainsi :

1°) \uline{0 fidélité} signifie que pour \(\forall\) hom \(u:\Gamma\to A\), on a \(\pi_1(\unsure{A^u})=0\)

\inserted{\uline{NB} \(A^u\) et \unsure{\(A^{0u}\)} ont \(\overline{\mathrm m}\) \unsure{composante neutre|composantes neutres}, donc la condition s'écrit : \fbox{\(\pi_1(\unsure{(A^0)^u})=0\)}.}

2°) \uline{1-fidélité} signifie que de plus, si \(u,v\) \deleted{sont deux} hom \(\Gamma\rightrightarrows A\), et \deleted{\unsure{\(\alpha,\beta\in\mathrm{Transp}\)} \illegible{1}} \(\alpha,\beta\in\mathrm{Transp}_A(u,v)\) \deleted{\illegible{1} par \illegible{2}} \inserted{ont \(\overline{\mathrm m}\) image dans \(\mathrm{Transp}_{\pi_0}(\uline{u},\uline{v})\)} \deleted{alors \illegible{1} \(\alpha,\beta\) \illegible{3} \(A\to\pi_0\) sont conjugués, donc \unsure{\((u,v)\)} \illegible{1}} \deleted{\illegible{1}} alors \(\alpha,\beta\) sont dans une \(\overline{\mathrm m}\) \deleted{\illegible{1}, \(\overline{\mathrm m}\) \unsure{identiques}), alors \(u\) et \(v\) sont} comp. connexe de \(\mathrm{Transp}_A(u,v)\). \deleted{\unsure{Cela}} \deleted{\unsure{conjugués entre eux}} \deleted{\illegible{1}} Pour le voir, OPS \(u=v\), et on est ramené à exprimer que l'application \(\pi_0(A^u)\to\pi_0\) est injective i.e. que \deleted{\unsure{\(A^u\)}} son noyau est réduit à \(1\), i.e. que \inserted{le sous-groupe ouvert \(\unsure{(A^u)\cap A^0=(A^0)^u}\)} de \(A^u\) est connexe, i.e
\margin{\illegible{1}}
\[
\fbox{\(\pi_0(\unsure{(A^0)^u})=0\)}
\]

\margin{278}

\scanpage{280}
\mspage{141}

3°) \uline{2-fidélité} signifie qu'\unsure{en} plus des conditions précédentes, qui assurent que pour \(u,v\) fixés, le foncteur \(\uline{\mathrm{Hom}}(u,v)\to\uline{\mathrm{Hom}}(\uline{u},\uline{v})\) est pl. fidèle, que celui-ci est ess. surjectif, i.e. que \unsure{si} \(u,v\) sont tels que \(\uline{u},\uline{v}:\Gamma\to\pi_0\) sont conjugués, \deleted{alors} (\unsure{si on veut égaux}) \(u,v\) sont déjà conjugués \inserted{\unsure{(par un él\(^{\mathrm t}\) de \(\uline{\pi_0}\))}}. En d'autres termes : pour \(\forall\) hom \(\uline{u}:\Gamma\to\pi_0=\unsure{A/A^0}\), s'il existe un relèvement de \(u\) en \(\Gamma\to A\), celui-ci est unique à conjugaison près.

\margin{\unsure{Il faut le dire} \illegible{1} \unsure{de façon} \unsure{plus forte} : \(\mathrm{Transp}(u,v)\to\mathrm{Transp}(\uline{u},\uline{v})\) \unsure{est} surjectif --- cela équivaut à dire que si \unsure{\(\uline{u}=\uline{v}\)}, alors \(u\) et \(v\) sont conjugués par un \uline{élément} \uline{de \unsure{\(A^0\)}}.}

4°) \uline{3-fidélité} signifie qu'en plus, \unsure{tout} hom \(\uline{u}:\Gamma\to\pi_0\) se relève en \(u:\Gamma\to A\).

En résumé, \inserted{si \(\pi_1(\unsure{A^0})=0\),} la 3-fidélité signifie que pour \(\forall\) hom \(\uline{u}:\Gamma\to\pi_0=\unsure{A/A^0}\), il existe \deleted{un} relèvement \(u:\Gamma\to A\), unique à conjugaison près, et qu'alors \(\unsure{(A^0)^u}\) est connexe et simplement connexe,

\margin{279}
